\documentclass[12pt]{article}
\usepackage[utf8]{inputenc}
\usepackage{amsmath, amssymb, amsthm}
\usepackage{geometry}

\geometry{a4paper, margin=1in}

\newtheorem{theorem}{Theorem}
\newtheorem{lemma}[theorem]{Lemma}
\newtheorem{corollary}[theorem]{Corollary}
\newtheorem{proposition}[theorem]{Proposition}
\theoremstyle{definition}
\newtheorem{definition}[theorem]{Definition}
\theoremstyle{remark}
\newtheorem{remark}[theorem]{Remark}

\title{Project Euler Problem 15: Lattice Paths}
\author{}
\date{}

\begin{document}
\maketitle

\section{Problem Statement}

Count the number of monotone lattice paths from $(0,0)$ to $(20,20)$ using unit steps right $(+1,0)$ and down $(0,+1)$.

\section{Formal Development}

\begin{definition}
A \emph{monotone lattice path} from $(0,0)$ to $(m,n)$ is a sequence of $m+n$ unit steps, each $R = (+1,0)$ or $D = (0,+1)$, with exactly $m$ right-steps and $n$ down-steps.
\end{definition}

\begin{definition}
Let $P(i,j)$ denote the number of monotone lattice paths from $(0,0)$ to $(i,j)$.
\end{definition}

\begin{theorem}[Closed Form]
\label{thm:closed}
$P(m,n) = \binom{m+n}{m}$.
\end{theorem}

\begin{proof}
A path is determined by choosing which $m$ of the $m+n$ steps are rightward. There are $\binom{m+n}{m}$ such choices.
\end{proof}

\begin{theorem}[Pascal Recurrence]
\label{thm:pascal}
$P(i,0) = P(0,j) = 1$ for all $i, j \ge 0$, and $P(i,j) = P(i-1,j) + P(i,j-1)$ for $i, j \ge 1$.
\end{theorem}

\begin{proof}
The boundary conditions hold since there is a unique all-$R$ path to $(i,0)$ and a unique all-$D$ path to $(0,j)$. For the recurrence: every path to $(i,j)$ arrives via its last step from either $(i-1,j)$ or $(i,j-1)$, and these cases are disjoint and exhaustive.
\end{proof}

\begin{lemma}[Consistency]
\label{lem:consist}
$P(i,j) = \binom{i+j}{i}$ satisfies Theorem~\ref{thm:pascal}.
\end{lemma}

\begin{proof}
By induction on $i+j$. Base cases are immediate. The inductive step follows from Pascal's identity: $\binom{i+j-1}{i-1} + \binom{i+j-1}{i} = \binom{i+j}{i}$.
\end{proof}

\begin{theorem}[Telescoping Product]
\label{thm:prod}
\[
\binom{2n}{n} = \prod_{k=1}^{n} \frac{n+k}{k}.
\]
Each partial product $\prod_{k=1}^{j} \frac{n+k}{k} = \binom{n+j}{j} \in \mathbb{Z}$.
\end{theorem}

\begin{proof}
$\prod_{k=1}^{n} \frac{n+k}{k} = \frac{(n+1)(n+2)\cdots(2n)}{n!} = \frac{(2n)!}{n!\,n!}$. Integrality of partial products follows from $\binom{n+j}{j} \in \mathbb{Z}$.
\end{proof}

\begin{corollary}
The multiplicative formula computes $\binom{2n}{n}$ using only integer arithmetic: each intermediate division is exact.
\end{corollary}

\begin{proposition}[Complexity]
The multiplicative formula runs in $O(n)$ time and $O(1)$ auxiliary space (with machine-word arithmetic for $n = 20$, since $\binom{40}{20} < 2^{63}$).
\end{proposition}

\begin{proof}
The loop has $n$ iterations. For $n = 20$, $\binom{40}{20} \approx 1.38 \times 10^{11} < 2^{63} \approx 9.2 \times 10^{18}$, so all intermediate values fit in a 64-bit integer.
\end{proof}

\section{Computation}

\[
\binom{40}{20} = \prod_{k=1}^{20} \frac{20+k}{k} = \frac{21}{1} \cdot \frac{22}{2} \cdot \frac{23}{3} \cdots \frac{40}{20} = 137{,}846{,}528{,}820.
\]

\section{Answer}

\[
\boxed{137846528820}
\]

\end{document}
